Nuclear power plants continuously acquire neutron and instrument-channel measurements for surveillance and condition monitoring. These records can reveal changes in detector behavior, signal integrity, and operating regimes, but are often sparsely labeled and difficult to align with a single sample-level fault taxonomy. Technical interpretation therefore continues to rely on authoritative guidance for online monitoring (OLM) and instrument-channel assessment~\citep{iaea2008olm1,iaea2008olm2,hashemian2011online}. Neutron-noise analysis further provides a non-intrusive route to studying reactor and instrumentation phenomena~\citep{pazsit2010noise,demaziere2022cortex}. Although data-driven diagnostic methods can be effective for predefined scenarios, their assumptions of fixed feature schemas and labeled fault classes are not always met by plant archives~\citep{tambouratzis2020neutron,elbordany2024fault}.

Large language models (LLMs) can help organize heterogeneous observations and technical knowledge, but free-form generation is unsuitable as a decision authority in a safety-relevant setting. Retrieval-augmented generation (RAG) conditions an answer on external evidence~\citep{lewis2020rag}; in nuclear applications, its value depends on whether the retrieved manuals genuinely support the statement being made. Recent nuclear RAG studies emphasize factuality, provenance, and failure modes such as incorrect retrieval and unsupported interpretation~\citep{reeves2024reliability,lin2025cws,mapes2025uprate,ndum2026radiant}. A practical open question is how to connect measurement-side findings to that evidence layer without asking an LLM to infer a root cause directly from a raw CSV file.

This work addresses that question through an evidence-constrained diagnostic-assistance workflow. A tool-using agent observes raw neutron-current series with minimal diagnostic priors and writes an auditable condition report. The report is then compressed into a structured alert summary that drives manual retrieval and citation-bearing explanation. The intended output is not an automatically confirmed fault diagnosis: it distinguishes observed data phenomena from manual-supported candidate interpretations, both of which remain subject to engineering review.

The reliability of this workflow depends critically on retrieval. A manual library may need to include both a small, pre-designated authoritative set and a larger set of same-domain supplements. Adding relevant supplements can improve apparent nearest-neighbor similarity while reducing the priority of the authoritative sources specified by a deployment policy. This tension is especially important for bilingual technical terminology and for overlap-based chunking, where a useful passage may not have the identical chunk identifier used as a query probe.

The contributions are threefold. First, we develop a zero-prior, tool-using workflow with a persistent workspace contract for converting neutron-current CSV data into auditable observations. Second, we introduce report-driven multi-stage RAG that maps structured observations to citable manual evidence. Third, we evaluate evidence retrieval under controlled core--supplement expansion, query language, chunking, and strict versus document-level relevance. The resulting evidence is used to characterize deployment boundaries rather than to claim autonomous fault identification. \Cref{fig:overview} summarizes the workflow.

\begin{figure}[t]
  \centering
  \includegraphics[width=0.98\linewidth]{fig1_overview}
  \caption{Evidence-constrained workflow. Deterministic probes turn a neutron-current CSV into an auditable condition report; a structured alert summary retrieves manual evidence for a citation-bearing diagnostic memo. Engineering review remains outside the automated chain.}
  \label{fig:overview}
\end{figure}
